% ============================================================================
\section{Limitations and Future Work}
\label{sec:limitations}
% ============================================================================

\subsection{Known limitations}
\begin{itemize}
  \item \textbf{One board simulates three terminals} (\cref{sec:one-board}). The terminals are
        isolated and talk only through the bus, but a separate-board transport (ESP-NOW or Wi-Fi)
        has not been implemented.
  \item \textbf{RFID is simulated} with buttons, because Wokwi has no RFID reader part; only
        \file{rfid.c} would change for a real reader.
  \item \textbf{The device does not send events to the dashboard yet.} The dashboard and its
        ingest API are tested on their own, with sample events.
  \item \textbf{Truncation of the newest log lines} cannot be detected by the hash chain on its
        own; an off-device copy (the dashboard) would cover that.
  \item \textbf{A few edge cases are not yet handled}, found while reviewing the code:
        \begin{itemize}
          \item pressing \code{A} during an admin session silently starts a new login;
          \item a card tapped after the login card step timed out shows ``Not an admin card'';
          \item a full occupancy ledger (64 people) does not record the next entry.
        \end{itemize}
\end{itemize}

\subsection{Future work}
\begin{itemize}
  \item Split the terminals onto three ESP32 boards using ESP-NOW for the same bus messages.
  \item A device-side sync task: Wi-Fi, a signed batch upload, and an upload cursor kept on the SD
        card so nothing is lost while offline.
  \item A local web page served by the hub over Wi-Fi.
  \item A closing-time warning on the LCD and buzzer 15 minutes before closing (the function
        \code{schedule_minutes_to_close()} is already written and tested).
  \item A real RFID reader (MFRC522) behind the existing \file{rfid.h} interface.
\end{itemize}
